%
% Settings file
%
%=======================================================================
%-------------This file controls various document settings--------------
%=======================================================================

%
% Should we compile for the Iliad (ebook reader), or regular A4
%
\newif\ifcompileiliad 
	\compileiliadfalse %Iliad (true) or regular A4 (false)?


%
% Set the fontsize and baselineskip here, use 13/15 for final document downscaling to 80% (UiB standard)
%
\newcommand{\TextSize}{13}
\newcommand{\BaseLineSkip}{15}
\newif\ifDownscaledFinalDoc
	\DownscaledFinalDoctrue		% Uib 13pt (true) or Regular 12pt (false)


%
% Compile draft (double line spacing)?
%
\newif\ifDraft
	\Draftfalse		% Draft (true) or Final (false)
	% Document settings file

%
% We use the book class
%
\ifcompileiliad
	\documentclass[10pt]{book}
\else
	\documentclass[12pt]{book}
\fi

%
% The geometry package allows consistent control of page layout
%
\ifcompileiliad
	\setlength{\parindent}{0pt}
	\setlength{\parskip}{3ex}
	\usepackage[papersize={122mm,160mm},
				hmargin=0.2cm,
				vmargin=0.2cm,
				tmargin=1.0cm,
				bmargin=0.4cm,
				dvips]
				{geometry}
\else
	\usepackage[paper = a4paper,         %To be shrunk to 80% when printed
				twoside,                 %Two-side mode, switches margins on 
				bindingoffset = 2mm,     %Offset for binding side of page
				hmargin = 25mm,          %Left and right margin
				vmargin = 25mm,          %Top and bottom margin
				dvipdfm]
				{geometry}
\fi	


% MY PACKAGES (MI)
\usepackage{subcaption}

%
% Font settings
%
\usepackage[T1]{fontenc}	% Activate Type 1 fonts
\usepackage{mathptmx}       % Use Times font, also for math
\usepackage[scaled]{helvet} % for sans serif fonts (\textsf{...} or \sffamiliy)
\usepackage{sectsty}        % Change section and chapter header
\allsectionsfont{\usefont{OT1}{phv}{bc}{n}\selectfont} % Set chapter/section header to narrow helvectica (arial-like)
%\usepackage[scaled]{luximono} % for monospaced fonts (\texttt{...} or \ttfamily)


%
% Packages we need
%
\usepackage{graphicx}       % Handles figures
\usepackage[utf8]{inputenc} % We want æøå
\usepackage{amsmath}
\usepackage{amsfonts}
\usepackage[mathcal]{euscript} % For caligraphy fonts
\usepackage{booktabs}       % Publication quality tables
\usepackage{setspace}       % Easy setting of line spacing
\usepackage{forloop}		% For-loops!
\usepackage[numbers,
			sort&compress
			]{natbib}       % Sort numerical keys for multiple cites
\usepackage{hypernat}       % Avoid breaking 'backref' option of hyperref package
\usepackage{float}
\usepackage[dvipsnames]{color}   % named colors
\usepackage[procnames]{listings} % beautiful listings
\usepackage{units}				 % semantically represent numbers with units
     

%
% The hyperref package allows advanced hyperlinking functionality
%
\usepackage[dvipdfmx,        % Use dvipdf driver
			backref,         % List citing occurences in the References
			colorlinks,      % Colored links
			citecolor=blue,  % Color of cite links
			linkcolor=blue,  % Color of links
			urlcolor=blue,   % Color of urls
			]{hyperref}
%\newcommand{\aref}[1]{\autoref*{#1}} % Prevent links
\newcommand{\aref}[1]{\autoref{#1}} % Enable links


%
% Set fancy page header using fancyhdr package
% 
\usepackage{fancyhdr}
\pagestyle{fancy}

% Ensure that the chapter and section headings are in lowercase
\setlength{\headheight}{15pt}
\renewcommand{\chaptermark}[1]{\markboth{#1}{}}
\renewcommand{\sectionmark}[1]{\markright{\thesection\ #1}}

% Delete current section for header/footer
\fancyhf{}

% Define header/footer layout
\fancyhead[LE,RO]{\bfseries\thepage}
\fancyhead[LO]{\bfseries\rightmark}
\fancyhead[RE]{\bfseries\leftmark}
\renewcommand{\headrulewidth}{0.5pt}

% make space for the rule
\fancypagestyle{plain}{
	\fancyhead{} %get rid of the headers on plain pages
	\renewcommand{\headrulewidth}{0pt} % and the line
}

%
% Set header height to 30pt to make room for two rows in the header
% (useful if you have very long chapter names)
%
%\headheight 30pt


%
% Include user-defined macros
%
%====================================================
%-------------Macro definitions go here--------------
%====================================================

%
% Differentials
%
\newcommand{\tdiff}[2]{\ensuremath{\frac{d#2}{d{#1}}}}
\newcommand{\tdifforder}[3]{\ensuremath{\frac{d^{#2}#3}{d{#1}^{#2}}}}
\newcommand{\pdiff}[2]{\ensuremath{\frac{\partial#2 }{\partial#1}}}
\newcommand{\pdifforder}[3]{\ensuremath{\frac{\partial^{#2}#3}{\partial{#1}^{#2}}}}

%
% Linear algebra
%
\renewcommand{\vec}[1]{\ensuremath{\mathbf{#1}}}
\newcommand{\mat}[1]{\ensuremath{\mathbf{#1}}}
\newcommand{\tildemat}[1]{\ensuremath{\widetilde{\mat{#1}}}}


%
% Misc macros
%
\newcommand{\eref}[1]{~(\ref{#1})}
\renewcommand{\equiv}[0]{\ensuremath{:=}}
\newcommand{\etal}{\textit{et al. }}
\newcommand{\paperheader}[2]{\noindent\textbf{Paper #1}: \textit{#2}\\}
\newcommand{\paperitem}[3]{\noindent\textbf{Paper #1}: \textit{#2}\vspace{1em}\\\noindent #3\vspace{2em}}
\newcommand{\tfinal}{\ensuremath{T_{\text{f}}}}
\newcommand{\papernum}[1]{\textbf{#1}}


%
% Theorem environments
%
\newtheorem{definition}{Definition}


%
% Counters
%
\newcounter{ct}

%
% For loop to include papers
%
\newcommand{\includepaperpages}[2]
{
	\forloop{ct}{0}{\value{ct} < #2}
	{
		\begin{figure}[ht]
			\includegraphics[width=.99\textwidth]{#1\the\value{ct}.ps}
		\end{figure}
		\clearpage
	}
}


%
% TeX is very proud of its hyphenation engine, to the point where it 
% hyphenates _everything_ just to show off.  Increasing penalty for 
% hyphenation, and increase tolerance for underfull boxes can make 
% the text easier to read, at the expense of making the text less aligned
% to the right margin
%
\hyphenpenalty=10
\tolerance=1000


% Redefine include command -> input command
%\renewcommand{\include}{\input}

% Include only these chapters
%\includeonly{introduction}


%
% Set up the frontpage, with author, title, etc
%
\title{
{\fontsize{28}{30}\usefont{OT1}{phv}{bc}{n}\selectfont Optimization Problems with Communication}
	\author{
	\textbf{Marika Ivanova}\vspace{3cm}\\
		\includegraphics[width=74mm]{figurer/uglo}\vspace{1em}\\
		Dissertation for the degree of Philosophiae Doctor (PhD)\vspace{3.5em}\\
		Department of Informatics\\
		University of Bergen
	}
	\date{Month year}
}

%====================================================
%------------------ BEGIN DOCUMENT ------------------
%====================================================
\begin{document}
%Cause all references in bibtex file to appear in the 'References' section, even
%if they are not explicitly cite'ed in the document
%\nocite{*}

%Set the fontsize and baselineskip, if something other than 10 or 12 is required
\ifDownscaledFinalDoc
	\fontsize{\TextSize}{\BaseLineSkip}
	\selectfont
\fi

%Double line spacing for draft
\ifDraft
	\doublespacing
\fi


%--------------------------------------------------------------------------------------------------
% FRONT MATTER
%--------------------------------------------------------------------------------------------------
\maketitle
\frontmatter
%\chapter{Scientific environment}

This is where you write the name of the subject/professional milieu involved in the study, faculty(faculties)/institute(s)/centre(s)/research groups/school of research.

Contractual joint venture partners should also be mentioned here. If these are represented by graphic symbols (logo etc.) the it should be on this page.

NB! Most organisations have guidelines concerning the use of their logo. These guidelines regulate who can use the logo, when the logo can be used and how. This is quite relevant when a logo is used together with information connected to other institutions/organisations. (A typical example would be rules concerning space between other written or graphic content.)

Remember to find out which guidelines apply and ensure that they are followed!
Most organisations will be able to provide a digital version of their logo in a suitable format.

\chapter{Acknowledgements}

I would like to thank\ldots

\chapter{Abstract}

%
This dissertation is a compilation of six research papers that are focused on three different topics summarized in the text.

The first three papers concern NP-hard problems arising in ad-hoc wireless communication discussed in Chapter 2.
In general, the task is to broadcast a message in a given network of wireless devices while minimizing the power consumption.
Problems in this category differ in requirements on the network connectivity, models of power consumption, and the ability of devices to initiate a signal transmission.
Some of the common features of these problems are that a device can simultaneously transmit a signal to all devices within its communication vicinity,
and that a signal can travel from its originator to its recipient via multiple intermediate devices.
The wireless networks are modeled by means of graph theory.
Solution techniques for these problems involve mainly methods of integer linear programming and inexact algorithms with or without performance guarantee.

The next paper is focused on the problem of minimum broadcast time.
Unlike the previous topic, devices in this problem are supposed to send a signal to up to one neighbouring device at a time.
The objective is to determine a sequence of signal transmission from a given set of source devices to the remaining ones while minimizing the time needed for spreading the signal.
Chapter 3 describes this problem in detail along with several related problems.
This problem is also studied from the integer linear programming as well as inexact algorithm perspective.
Continuous relaxations of the ILP models help to evaluate quality of the studied inexact methods.
The stronger model, the more accurate assessment it provides.

The last two papers are dedicated to problems belonging to path planning for multiple robots discussed in Chapter 4.
In general, these problems involve a group of agents (robots) initially deployed in an environment, and the task is to find a sequence of their moves
so that they reach pre-defined destination locations while optimizing a given criterion such as minimum makespan or minimum total arrival time.
The agents' movement must obey a set of given rules.
An extension of the problem considers agents divided into two (or more) adversarial teams, where the teams have either symmetric or asymmetric objectives.
After introducing the adversarial element, the problem of finding a winning strategy for a given team becomes PSPACE-hard, like many two player games with alternating turns.
%Additional requirement of preserving a communication possibility during the agets' movement has also been studied.

%The remainder of this chapter is structured as follows: 
%Section~\ref{sec:back} contains an overview of general underlying mathematical and algorithmical tools and methods employed in the attached papers.
%In Section~\ref{sec:smt}, we discuss ad-hoc wireless networks and describe several optimization problems in this area.
%In particular, we focus on the shared multicast tree problem and provide an insigt into its relation to other similar problems.
%Section~\ref{sec:mbt} introduces the minimum broadcast tree problem and discusses other relevant problems.
%In Section~\ref{sec:app}, we describe the problems in multi-robot path planning, and summarize their properties.
%An emphasis is given to the area protection problem.

%This is the introduction~\cite{deb01, pyprop}\ldots

%This is where you write the abstract, ref: 
%
%\vspace{1em}
%\noindent \textbf{Ӥ 6.6 Abstract of dissertation, press release}
%
%\vspace{1em}
%\noindent \textit{An abstract of the dissertation shall be prepared in English
%(1-3 pages), so that research environments in Norway and abroad may familiarize
%themselves with the results. The abstract shall accompany the dissertation.”}
%
%
%\vspace{1em}
%\noindent From  ”The Regulations for the degree philosophiae doctor (PhD) at the
%University of Bergen” 
%
%\vspace{1em}
%\noindent (There is an abstract template that has room for information about the debate,
%opponent etc that can be used as a handout.)

\chapter{List of papers}

\begin{enumerate}

\item Author, \textit{Sometitle}, Somejournal \textbf{1}, 1, 1905.

\end{enumerate}

\tableofcontents
\listoffigures
\lstlistoflistings

%--------------------------------------------------------------------------------------------------
% MAIN MATTER
%--------------------------------------------------------------------------------------------------
\mainmatter

%
% Main chapters
%

%
% CUSTOM THEOREMS (MI)
%
\newtheorem{theorem}{\textbf{Theorem}}
%\newtheorem{observation}[theorem]{\textbf{Observation}}
%\newtheorem{definition}{\textbf{Definition}}
\newtheorem{problem}[theorem]{\textbf{Problem}}
%\newtheorem{proposition}[theorem]{\textbf{Proposition}}
\newtheorem{remark}[theorem]{\textbf{Remark}}
%\newtheorem{lemma}[theorem]{\textbf{Lemma}}


\chapter{Introduction}
%
This dissertation focuses three optimization problems of which common denominator is communication among multiple devices.

The first two topics concern NP-hard combinatorial network optimization problems, and can naturally be formulated as integer linear programs (ILP).
Instances of limited sizes can be solved to optimality using traditional methdos for solving ILP such as Branch and Bound (B\&B).
In case of larger instances, where finding optimal solution is prohibitively time consuming, methods without guarantee of optimality such as heuristic algorithms are employed.
ILP models are also useful whenever the optimality has to be given up due to the excessive instance size.
Continuous relaxations of ILP models help to evaluate quality of a studied inexact method.
The stronger model, the better assessment it provides.
Devising strong ILP models is therefore of our great interest.
%A common optimization criterium in is energy consumption, as the devices are typically heavily energy constrained due to the use batteries as energy sources.
%Another criterium is time during which a communication session takes place.
The cornerstone of the first two problems is a set of devices of fixed locations with the ability of communicating with each other.
The task is always to disseminate a message initiated at certain source devices to the remaining ones while abiding by certain communication rules and restrictions.
In practical applications, there are various optimization criteria in network communication problems.
The first two problems studied herein require minimizing the power consumption and time during which a communication session takes place.

The last problem belongs to the artificial intelligence (AI) world, particularly to the family of multi-robot path planning problems.
The most general variant, cooperative path-finding \cite{silver05} is NP-hard, and assumse a set of agents deployed in an environment.
Each agent is assigned a target location in the environment, and their task is to reach these locations by a simultaneous movement while obeying several rules.
There are three most common optimization criteria \cite{yu13} to be pursued:
\begin{itemize}
	\item Minimum total arrival time - total number of time steps that the agents need before arrivin in their targets.
	\item Minimum makespan - the number of time steps needed by the latest arriving agent
	\item Minimum total distance - total number of moves performed by the agents
\end{itemize}
Various modifications have been proposed in the literature, and they will be discussed in detail in Section \ref{sec:app}.
An extension of the problem considers agents divided into two adversarial teams, where the teams can have either symmetric or asymmetric objectives.
After introducing the adversarial element, the problem becomes PSPACE-hard, like many two player games with alternating turns.
Additional requirement of preserving a communication possibility during the agets' movement has also been studied.

The remainder of this chapter is structured as follows: 
Section~\ref{sec:back} contains an overview of general underlying mathematical and algorithmical tools and methods employed in the attached papers.
In Section~\ref{sec:smt}, we discuss ad-hoc wireless networks and describe several optimization problems in this area.
In particular, we focus on the shared multicast tree problem and provide an insigt into its relation to other similar problems.
Section~\ref{sec:mbt} introduces the minimum broadcast tree problem and discusses other relevant problems.
In Section~\ref{sec:app}, we describe the problems in multi-robot path planning, and summarize their properties.
An emphasis is given to the area protection problem.

%This is the introduction~\cite{deb01, pyprop}\ldots

\section{Background}\label{sec:back}

%\subsection{Notation and Terminology}

\subsection{Graph Terminology}

A \emph{graph} $G$ is a pair $(V,E)$ of \emph{vertices} and \emph{edges}, where $E\subseteq {{V}\choose{2}}$.
If this inclusion is an equality, $G$ is said to be \emph{complete}.
A \emph{path} in graph is a sequence of edges connecting a sequence of distinct vertices.
A graph is said to be \emph{connected} if there exists a path between every two vertices, otherwise it is \emph{disconnected}.
A \emph{cycle} of a graph $G$ is a subset of $E$ that form a path such that the first node of the path corresponds to the last. 
If $G$ contains a cycle, $G$ is called, \emph{cyclic}, otherwise it is called \emph{acyclic}.
\begin{definition}
A tree is a graph that is connected and acyclic.
\end{definition}
For a vertex $v\in V$, a \emph{neighbourhood} of $v$ (open neighbourhood), denoted as $N(v)$, is the set of vertices adjacent to $v$.
The size of neighbourhood of $v$ is called a \emph{degree} of $v$.
A \emph{subgraph} of $G=(V,E)$ is a graph $G'=(V',E')$ such that $V'\subseteq V$ and $E'\subseteq E$.
This relation is often written as $G\subseteq G'$.
\begin{definition}
A spanning tree of graph $G=(V,E)$ is a tree $T=(V_T,E_T)$ such that $T\subseteq G$ and $V_T=V$.
\end{definition}

In a \emph{weighted graph} $G=(V,E,w)$, a \emph{weight} or \emph{cost} $w_e$ is associated witch each edge $e\in E$.
A weight of $G$ is $\sum_{e\in E}w(e)$.
A spanning tree of $G$ with minimum weight is called a \emph{minimum spanning tree} of $G$.
Similar concept is used in paths in graphs.
A \emph{shoretst path} from $u$ to $v$ in a weighted graph is a path consisting of edges of minimum sum of weights connecting $u$ and $v$.
\begin{definition}
	For a graph $G=(V,E)$ and a subset of vertices $D\subseteq V$, a Steiner tree of $G$ and $D$ is a tree $T=(V',E')$ such that $T\subseteq G$ and $D\subseteq V'$.
\end{definition}
Analogously in weighted graphs, a \emph{minimum Steiner tree} is a Steiner tree of minimum weight.

If edges have a direction associated with them, we call such a graph a \emph{directed graph}, and its edges are referred to as \emph{arcs}
Let $G=(V,A)$ be a directed graph. 
By \emph{in-degree} and \emph{out-degree} of a vertex $v\in V$ we mean the sizes of sets $|\{(u,v): u\in V, (u,v) \in A\}|$ and $|\{(v,u): u\in V, (v,u) \in A\}|$ , respectively.
An \emph{arborescence} rooted at vertex $r$ is a directed tree with arcs directed from $r$.

\subsection{Combinatorial Optimization}

Combinatorial optimization (CO) is a part of applied mathematics that taclkes optimization problems over discrete structures.
It combines methods from graph theory, linear programming, combinatorics and the theory of algorithms.
In this section, we briefly introduce main concepts in CO used later in the text.
For a comprehensive rendition of this topic, an interested reader is referred to \cite{wolsey98} and \cite{nemhauser88}.

Combinatorial problems arise in many areas of computer science with a wide range of applications in various industrial disciplines 
such as production scheduling, logistics, communication network design and many more.
The core solving a problem by methods of CO is the identification of a discrete mathematical structure hidden in the problem,
and finding a sufficient abstraction.

CO concerns problems of minimization or maximization of an \emph{objective function} of several variables 
subject to inequality and equality \emph{constraints} and integrality restriction on at least some of the variables.
In this work, both the objective function and constraints are assumed to be linear.
Combinatorial problems are often formulated as \emph{mixed integer linear programs} (MILP) of the standard form
\begin{equation}
\begin{array}{r@{}l}
	\max\limits_{x,y}c^\top x + h^\top y & \\
	\text{subject to}& \\
	  Ax + Gy &\leq b, \\
	  x \in \mathbb{Z}^{n}_+, y & \in\mathbb{R}^{p}_+. 
\end{array}
\end{equation}

The problem instance is specified by the input data $c\in \mathbb{R}^n$, $h \in \mathbb{R}^n$, $A \in \mathbb{R}^{m\times n}$  $G \in\mathbb{R}^{m\times p}$ and $b \in \mathbb{R}^m$.
A MILP that is not in the standard form, for example if the objective is to minimize or if the constraints contain equalities, can be straightforwardly converted into the standard form.
If the integrality constraints are not present, we talk about a \emph{linear program} (LP)


The set of points $S=\{(x_0,y_0):  x_0 \in \mathbb{Z}^{n}_+, y_0  \in\mathbb{R}^{p}_+, Ax_0 + Gy_0 \leq b\}$ is called the \emph{feasible region},  
and a point $(x_0,y_0)\in S$ is referred to as a \emph{feasible point} (feasible solution) with \emph{objective function value} $c^\top x_0 + h^\top y_0$. 
A feasible point $(x^*,y^*)$ is called an \emph{optimal solution} if for every feasible points $(x_0,y_0)$ we have that $c^\top x_0 + h^\top y_0 \leq c^\top x^* + h^\top y^*$. 
Expression $c^\top x^* + h^\top y^*$ is then called the \emph{optimal value}. 

\subsubsection{Relaxation and Bounds}

\begin{definition}
	Let $S\subseteq \mathbb{R}^n$ and $\mathcal{F}$ be a MILP $\max\{f(x):x\in S\}$.
	The problem $\mathcal{R}: \max\{g(x):x\in T\}$ is a relaxation of $\mathcal{F}$ if and only if
	\begin{enumerate}
		\item $T\supseteq S$, and
		\item $g(x)\geq f(x) \forall x\in S$.
	\end{enumerate}
\end{definition}


Let $z^*$ and $\underline{z}$ be the optimal value of a MILP and its relaxation, respectively. 
Further, let $\bar{z}$ be the objective function value of some feasible point.
Then, $\underline{z}\leq z^* \leq \bar{z}$.
Values $\underline{z}$ and $\bar{z}$ are referred to as a lower bound and an upper bound, respectively.

A \emph{combinatorial relaxation} of a MILP is achieved by omitting one or more constraints. 
By omitting the integrality constraints of a MILP $\mathcal{F}$ we obtain its \emph{continuous relaxation}, also called \emph{LP relaxation}, denoted as LP$(\mathcal{F})$.

\subsubsection{Duality}

\subsubsection{Solving methods}

Linear programs are commonly solved in polynomial time by the \emph{simplex algorithm}.
A MILP can be solved by the \emph{branch and bound} (B\&B) method which may take exponential time.
These and other algorithms are an integral parts of most of the modern solvers such as CPLEX and GUROBI.

\subsubsection{Problems and Complexity}\label{sect:probcomp}

A \emph{computational problem} (problem) is an infinite collection of instances together with a solution for each instance.
A problem that can be posed as a yes-no question of the input values is referred to as \emph{decision problem}.
An example of a decision problem is the clique problem: Given a graph $G$ and an integer $k$, is there a clique in $G$ of size at least $k$?
In \emph{optimization problems}, the task is to find a ``best possible'' solution from among the set of all feasible solutions the problem instance.
The optimization version of the clique problem asks for finding the maximum clique in a given graph $G$.
An optimizatoin problem can be solved by answering a sequence of decision problems:
Assume there is an oracle that is able to answer the clique problem for a given $(G,k)$.
The maximum clique problem can then be solved by answering its optimization version for $k=1,2,\dots$ until the answer is ``no'' for some $k'$, and so the maximum clique has the size $k'-1$.

Throughout this thesis, we use several well known concepts from the complexity theory, which we state in the following.
Detailed explanations of the terminology can be found in any textbook on this topic such as \cite{sipser06}.

An \emph{algorithm} is a procedure that solves a given problem in a finite number of steps.
Computational complexity of an algorithm is the amount of time needed for its run, and is measured in terms of the input size.
A \emph{polynomial} algorithm runs in time $\mathcal{O}(n^c)$, for some constant $c$ and input $w$ of size $|w|=n$.
A \emph{verifier} is an algorithm that determines whether a given certificate is a proof to the fact that $w$ is a yes-instance.
An example of a certificate in the clique problem is some subset of nodes of size $k$.
It can be verified in polynomial time whether there exists an edge between every two nodes.
\begin{definition}
	$P$ is the class of decision problems for which there exists a polynomial algorithm that solves them.
\end{definition}
\begin{definition}
	$NP$ is the class of decision problems for which there exists a polynomial verifier. 
\end{definition}
\begin{definition}
	Problem $X$ is polynomial time reducible to problem $X'$, if a polynomial computable function $f$ exists where for every $w$, 
	$w$ is a yes-instance of $X\Leftrightarrow f(w)$ is a yes instance of $X'$.
\end{definition}
\begin{definition}\label{def:npc}
	Decision problem $X$ is NP-complete if it satisfies:
	\begin{enumerate}
		\item $X$ is in NP.
		\item Every $X'$ in NP is polynomial time reducible to $X$.
	\end{enumerate}
\end{definition}
Remark:
A verifier does not decide whether a certificate is an optimal solution to a given optimization problem instance. 
When addressing optimization problems, we consider only the second property in Def.~\ref{def:npc}. 
If this property is satisfied, we say that the problem is NP-hard.

There are many well know examples of NP-hard problems, and they can be formulated as ILP.
Therefore, ILP itself is also NP-hard.

\begin{definition}
	PSPACE is the class of decision problems for which there exists an algorithm that solves them in a polynomial space.
\end{definition}
\begin{definition}\label{def:psc}
	Decision problem $X$ is PSPACE-complete if it satisfies:
	\begin{enumerate}
		\item $X$ is in PSPACE.
		\item Every $X'$ in PSPACE is polynomial time reducible to $X$.
	\end{enumerate}
\end{definition}
If $X$ satisfies only the second condition in Def.~\ref{def:psc} we say that $X$ is \emph{PSPACE-hard}.


\subsection{Path Planning}

Basic path-finding in (weighted) graphs is a well known task in computer science. The aim 
is to find a route between two selected vertices. We also use the term \emph{source} and
\emph{target} to refer to the initial vertex and target vertex, respectively.

%\subsubsection{Dijkstra's algorithm}

A shortest path can be found using \emph{Dijkstra's algorithm} \cite{dijkstra59}, which guarantees optimality.
Dijkstra’s Algorithm visits vertices in the graph starting with the source. 
It then repeatedly inspects the vertex which was not yet visited that lies closest to the source. 
It expands outwards from the source until it reaches the target. 
A drawback of this approach is that it may expand too many vertices that later turn out to lie very far from the target.

%\subsubsection{Best first search}

Another possibility is the greedy \emph{best first search} algorithm which can find the target faster, but the selected path is not guaranteed to be optimal.
Instead of expanding vertices close to the source, it selects those close to the target.
Obviously, which node is exactly the closest is not known without a previous preprocessing, and so it uses a heuristic estimation which guides the way towards the target.

%\subsubsection{A*}

The \emph{A*} algorithm combines the advantages of Dijkstra's and best first search algorithm.
It guarantees to find an optimal path and also exploits a heuristic estimation of the distance to the target which helps to avoid expanding unsuitable vertices.

The heuristic estimation is particularly useful when there is an incomplete information about the graph.
For that reason, A* is popular in video game development, where the exact distances can not be predicted due to a frequent change of the environment. 

Let us denote the source and target by $s$ and $t$, respectively.
The following values are considered for every vertex $v$:
\begin{itemize}
	\item $g_v$, a distance from $s$ to $v$,
	\item $h_v$, heuristic estimation of distance between $v$ and $t$, 
	\item $f_v=g_v+h_v$.
\end{itemize}


\section{Power Minimizing Trees in Ad-hoc Wireless Netowrks}\label{sec:smt}

A Wireless Ad-hoc Network (WANET) is a decentralized type of wireless network.
Unlike cellular and wired networks, WANET does not depend on any pre-existing infrastructure.
The set of wireless communication devices with fixed locations are connected via wireless links depending on 
distance between the nodes, thier transmission power, error control scheme, background noise and interference.
Each device is equipped with an omnidirectional antenna with a communication range schematically illustrated in Fig. \ref{fig:omniantenna}.
Hence, a signal reaches all nodes within the communication range of its sender.
This range is determined by the power assigned to the sender, and this power can be changed over time.
The power necessary for relaying a signal to multiple devices is the maximum rather than the sum of the powers necessary to reach all intended receivers.
This feature is referred to as the \emph{wireless advantage} \cite{wieselthierXX}
Each device works as a transceiver, which means that it can both transmit and receive a signal.

The absence of infrastructure implies their quick deployment and simple configuration.
The recent years have withnessed an increasing interest in WANETs, motivated by numerous both civil and military applications and by the continual progression in wireless technologies.
Digital battlefield, disaster management, luggage handling in airports, context aware computing and mobile commerce are examples of the growing list of potential applications of WANETs \cite{younis06}.

, makes them useful in emergency situations such as disaster relief and military operations.

Fig. \ref{fig:communication} depicts different types of communication in wireless networks.
Consider a group of devices and a specific sender who initiates a signal transmission. 
A \emph{unicast} means that there is a single recipient.
Whenever the recepients are all the remaining devices within the group, we talk about a \emph{broadcast}.
Finally, a \emph{multicast} means that only some of the devices must receive the message. 
The remaining ones do not have to, but they can serve as intermediate devices forwarding the signal.

\subsection{Related Problems}

\section{Minimum Broadcast Trees}\label{sec:mbt}

\section{Multi-Agent Path Finidng}\label{sec:app}

In multi-robot path finding (MPF) we consider an environment with several identical moving entities referred to as 
\emph{agents}. Source and target locations are uniquely determined for each agent. The objective
is to find a route for every agent from its source to its target.
Agents must not collide with obstacles and other agents that are also moving along
planned routes towards their own targets.
The environment is modeled as an undirected graph, where the agents are placed in the vertices, and move along the edges from one vertex to another. 
Continuous time is divided into discrete time steps, where the relocation of an agent from one vertex to its neighbor takes exactly 1 time step.

Different variants and restrictions have been studied.

\emph{Muliti-agent path-finding} (MPF) \cite{source1}  supposes a group of agents in a given
environment, where initial and target positions are determined for each agent.
All individual agents must avoid collisions. Movement of the agents is carried
out in discrete time steps. Agent can shift from one vertex to its neighbor on
condition that the neighbor is either unoccupied or is being left by other agent
in the same time step. At most one agent is allowed to pass an edge within one
time step. That is, agents are not allowed to exchange their positions within one
time step.

\emph{Pebble motion on graphs} \cite{source1,source2} is a very similar problem as MPF. It can
be regarded as a restricted variant of MPF. The difference consists in rules for
movement. While MPF enables entering a vertex that is simultaneously being
left by other agent, such transfer is not permissible in pebble motion. As an
illustration we can mention 15 puzzle also known as Lloyd’s 15 \cite{lloydXX}.

\emph{Cooperative Path-finding} \cite{cpf} is a special case of MPF where each agent is
assumed to have a full knowledge of all other agents and their planned routes.
Precisely speaking, solving algorithm can take into account paths planned for
agents that were processed earlier, and adjust paths searched later according to
them.

\subsection{Scenarios with adversaries}

\subsubsection{Adversarial Cooperative Path Finding}

\subsubsection{Area Protection Problem}

\subsubsection{Area Protection Problem with Communication maintenance}

         % Introduction
\chapter{Contribution of the Thesis and Prospective Research}

This thesis is a compilation of six papers in which three independent topics are looked into.
The first three papers are focused on ad-hoc wireless networks discussed in Chapter~\ref{chap:wanet}.
The fourth paper deals with the broadcast time problem addressed in Chapter~\ref{chap:mbt}.
These four papers share the common characteristic that integer programming is an essential solution method.
In this sense, the last two papers are distinguished from the former.
They are devoted to adversarial variants of path finding for multiple robots summarized in Chapter~\ref{chap:app}, 
which are commonly approached by methods prevalent in the field of artificial intelligence.
%
% Paper I
%
\section{Paper I:  Shared Multicast Trees in Ad-hoc Wireless Networks}

The contribution of Paper I lies in introducing the Shared Multicast Tree (SMT) problem, a generalization of the Shared Broadcast Tree (SBT) problem.
An ILP model for SMT is proposed by modification of a model for SBT presented in \cite{yuan12}.
Additional constraints related to non-destination nodes are included in the model, and constraints in the existing model for SBT are quantified with respect to the non-destinations.
The presented model is subsequently proved to be a correct formulation of SMT.
Appendix A of Paper I contains details of the proof.

Further, several inexact construction methods are proposed.
The first two methods are based on construction of a solution to different problems, specifically MST and MEB, respectively. 
The construction is followed by local search, which identifies non-destinations whose presence in the solution is disadvantageous.
These non-destinations are then eliminated from the solution.
The third method is an algorithm devised specifically for SBT/SMT.
Its main idea is to gradually expand a solution by adding new edges, anticipate subtree sizes in the resulting solution, 
and thereby give an estimate of the final objective function value each time a new edge is appended. 

The experimental part is divided into two sections.
The first section focuses on investigation of the instance sizes that are solvable by the proposed ILP model using the CPLEX solver.
It is also studied how the increasing number of destinations and non-destinations affects the solution time and the objective function value.
In the second section, the comparison of the inexact algorithms indicates that the method based on subtree size anticipation outperforms the other two algorithms.

%
% Paper II
%
\section{Paper II: The Shared Broadcast Tree Problem and MST}

Because of the limited size of instances practically solvable to optimality, following from the computational complexity of SBT, 
approximability and approximation algorithms are often researched.
Paper II presents an instance of SBT for which the algorithm that construct a MST yields a solution to SBT with objective function value six times the optimum, 
proving that the approximation ratio of the algorithm that constructs a MST as a solution to SBT is at least 6.
In addition, experimental results then reveal that the ratio between the optimal objective function value,
and the objective function value of MST solutions in randomly generated instances is much  more favourable than the lower bound on the approximation ratio.

%
% Paper III
%
\section{Paper III: Integer Programming Formulations for the Shared Multicast Tree Problem}

Two ILP formulations for SMT are developed in Paper III.
The first model is based on broadcast trees, whereas the second one adapts several network flow techniques presented in \cite{polzin01}.
Both models are subsequently extended by redundant variables and corresponding constraints which strengthen the formulations.
The models are further strengthened by introducing valid inequalities.
A theoretical study of the models proves that network flow based models are at least as strong as corresponding models built on broadcast trees.
The experimental evaluation discovers instances showing that flow based models are in fact stronger.

Experimental evaluation further exhibits a profound trade-off between the time necessary for solving LP relaxation of the models, and the strength of the lower bound obtained.
The LP relaxation of the strongest model yields an integral solution in most of the instances.
However, its running time rules out its practical usability.
For addressing this issue, a constraint generation (CG) technique is developed, which increases the size of practically solvable instances.
Lower bounds are also yielded during the course of standard branch and bound algorithm.
A comparison of lower bound produced from CG and branch and bound shows that CG is able to obtain stronger lower bounds within the selected time limit of 20 minutes.

Paper III also contains an adaptation of a metaheuristic algorithm from \cite{pajor18}, originally developed for the minimum Steiner tree problem, 
combined with local search methods developed in Paper I.
The algorithm is able to solve a vast majority of tested instances to optimality. 
However, the optimality is proved only in instances for which it is possible apply branch and bound and solve them to optimality within a practical time.
In larger instances, the computed solutions are not proved to be optimal, but the results indicate a very good potential of the metaheuristic algorithm. 
%
% Paper IV
%

\section{Paper IV: Computing the Broadcast Time of a Graph}

We develop a straightforward ILP model for the Minimum Broadcast Time (MBT) problem, as the existing literature focusing on matematical models for MBT is rather scarce. 
Recently published \cite{chu17} contains a non-linear mathematical formulation of MBT, and authors in \cite{desousa18} present an ILP model.
The models in these works serve for the purpose of formal description of MBT, and are not investigated further as potential solution methods.
On the contrary, an exact method that exploits the problem characteristics and iteratively solves a decision version of the presented ILP model is devised in Paper IV.
This method is also applied to the LP relaxation of the model, and the outcome indicates that it yields strong lower bounds, often coinciding with the optimum.

Besides the continuous relaxation, analytical and combinatorial lower bounding techniques are studied. 
According to the experimental evaluation, these methods provide weaker lower bounds than the LP relaxation.
Upper bounds are obtained by a greedy algorithm of which the main feature is an iterative construction of broadcast trees of restricted size.
This algorithm is parametrized by the size limitation of the broadcast trees, and the larger trees are searched, the tighter upper bound can be achieved.

The numerical experiments also provide an insight into the relation between some of the graph properties and its broadcast time. 
It has been observed that graphs with more nodes as well as denser graphs tend to have its broadcast time closer to its trivial lower bound $\log(n)$ (see Sect.~\ref{sec:bg}).
Also, increasing size and density of the graph instances narrows the gap between upper and lower bounds.


%
% Paper V
%
\section{Paper V: Area Protection in Adversarial Path-finding Scenarios with Multiple Mobile Agents on Graphs}

Paper V introduces the Area Protection Problem (APP) as a modification of the previously studied Adversarial Cooperative Path Finding (ACPF) \cite{ivanova14} problem.
Its main contribution lies in the proof that APP is PSPACE-hard.
This result is achieved by demonstrating a polynomial-time reduction from \textsc{True Quantified Boolean Formula} (TQBF).

Several strategies are investigated for the team of defenders.
A building block of all the considered strategies is a so called single stage \emph{destination allocation}, in which a node is assigned to each defender at the beginning of the agents' movement.
The defenders then try to reach these destinations by any CPF algorithm, in this case we selected LRA* (see Sect.~\ref{sec:lrastar}).
A destination can be regarded as a target node for defenders, the difference is that while a target is part of the input, the destination is determined by a solution method.

The strategies differ in the approach of target allocation.
The two simplest methods, random and greedy allocation, select attackers' targets as destinations for defenders.
A more sophisticated method, \emph{bottleneck simulation}, runs a simulation of attackers' movement and tries to predict locations frequently passed by attackers.
These locations are assumed to be bottlenecks in the environment, and blocking them may prevent a larger number of attackers from reaching their targets.

The experiments are carried out on different environment types and different positions of teams.
This method is particularly successful in environments rich on bottlenecks, as it successfully identifies them.
%
% Paper VI
%
\section{Paper VI: Maintaining Ad-hoc Communication Network in Area Protection Scenarios with Adversarial Agents}

An extension of APP, in which defenders are required to maintain the possibility of communication between each other, is presented in Paper VI.
The possibility of communication is modeled by connectivity of the communication graph, whose nodes are the defenders' locations, 
and whose edges connect node pairs between which agents can communicate. 

We approach this problem by dividing the defenders into \emph{communicators} and \emph{occupiers} with different purposes. 
Occupiers have the same task as regular defenders, i.e., they intend to prevent attackers from reaching their targets, 
while communicators are supposed to ensure the connectivity maintenance by moving to suitable nodes.

From a theoretical point of view, we show that the problem whether it is possible for communicators to maintain communication when all occupiers reach their destinations is NP-complete, 
which was proved by reducing \textsc{Vertex Cover} to it.
  % Introduction to the papers
\chapter{Scientific results}


%--------------------------------------------------------------------------------------------------
% Paper I
%--------------------------------------------------------------------------------------------------
\chapter*{Paper I}
\section{Shared Multicast Trees in Ad Hoc Wireless Networks}

\noindent Marika Ivanova\\

%\noindent \textit{Physical Review} B, \textbf{76}, 035303 (2007)
\noindent In: Cerulli R., Fujishige S., Mahjoub A. (eds) Combinatorial Optimization. ISCO 2016. Lecture Notes in Computer Science, vol 9849. Springer, Cham
\cleardoublepage

% This will include the bursted pdf. Remember to update the number of pages
% (last parameter) to the actual number of pages
\includepaperpages{../papers/paper-1/page_}{14}

%--------------------------------------------------------------------------------------------------
% Paper II
%--------------------------------------------------------------------------------------------------
\chapter*{Paper II}
\section{The Shared Broadcast Tree Problem and MST}

\noindent Marika Ivanova\\

%\noindent \textit{Physical Review} B, \textbf{76}, 035303 (2007)
\noindent Electronic Notes in Discrete Mathematics, vol 55, 2016
\cleardoublepage

\includepaperpages{../papers/paper-2/page_}{4}
%--------------------------------------------------------------------------------------------------
% Paper II
%--------------------------------------------------------------------------------------------------
\chapter*{Paper III}
\section{Integer Programming Formulations for the Shared Multicast Tree Problem}

\noindent Marika Ivanova, Dag Haugland\\

%\noindent \textit{Physical Review} B, \textbf{76}, 035303 (2007)
\noindent Under revision in Journal in Combinatorial Optimization
\cleardoublepage

\includepaperpages{../papers/paper-3/page_}{31}
               % The papers, with cover pages

%
% Appendices
%
%\appendix
%\chapter{List of NP-hard Problems}


\begin{problemA}
  \problemtitle{\textsc{3D-Matching}}
  \problemdefinition{
    Let $X$, $Y$, and $Z$ be finite, disjoint sets, and let $T$ be a subset of $X\times Y \times Z$. 
    The set $M \subseteq T$ is a 3-dimensional matching if for any two distinct triples $(x_1, y_1, z_1) \in M$ and $(x_2, y_2, z_2) \in M$, we have $x_1 \neq x_2, y_1 \neq y_2$, and $z_1 \neq  z_2$.  
  }
  \probleminput{A set $T$ and an integer $k$.}
  \problemquestion{Does there exist a 3-dimensional matching $M \subseteq T$ with $|M| \geq k$?}
\end{problemA}

\begin{problemA}
  \problemdefinition{A boolean formula is in 3CNF if it is a conjunction of clauses, where each clause is a disjunction of at most 3 literals}
  \problemtitle{\textsc{3-SAT}}
  \probleminput{A Boolean formula $varphi$ in 3CNF.}
  \problemquestion{Is there a truth assignment to variables of $\varphi$ such that $\varphi$ evaluates to TRUE?}
\end{problemA}

\begin{problemA}
  \problemdefinition{See \textsc{3-SAT}. }
  \problemtitle{\textsc{3-3-SAT}}
  \probleminput{A Boolean formula $varphi$ in 3CNF with an additional requirement that each variable appears at most three times and each literal at most twice.}
  \problemquestion{Is there a truth assignment to variables of $\varphi$ such that $\varphi$ evaluates to TRUE?}
\end{problemA}

\begin{problemA}
  \problemtitle{\textsc{At Least 1-in-3-SAT}}
\end{problemA}

\begin{problemA}
	\problemtitle{\textsc{Boolean Satisfiability (SAT)}}
  \probleminput{A Boolean formula $varphi$.}
  \problemquestion{Is there a truth assignment to variables of $\varphi$ such that $\varphi$ evaluates to TRUE?}
\end{problemA}

\begin{problemA}
  \problemtitle{\textsc{Clique}}
  \problemdefinition{A clique in a graph $G$ is a complete subgraph of $G$.}
  \probleminput{A graph $G$ and a an integer $k$.}
  \problemquestion{Does there exist a clique of size at least $k$ in $G$?}
\end{problemA}

\begin{problemA}
	\problemtitle{\textsc{Minimum Broadcast Time (MBT)}}
  \problemdefinition{
	 Let $G=(V, E)$ be a graph and let $S\subseteq V$ be a subset of vertices, called terminals. 
	A Steiner tree in $G$ is a tree that spans $S$.		
  }
  \probleminput{A graph $G=(V,E)$, non-negative edge-weights $c$, and $S\subseteq V$.}
\problemquestion{Find a minimum-weight Steiner tree in $G$.}
\end{problemA}

\begin{problemA}
	\problemtitle{\textsc{Minimum Energy Broadcast (MEB)}}
  \problemdefinition{
	 Let $G=(V, E)$ be a graph and let $S\subseteq V$ be a subset of vertices, called terminals. 
	A Steiner tree in $G$ is a tree that spans $S$.		
  }
  \probleminput{A graph $G=(V,E)$, non-negative edge-weights $c$, and $S\subseteq V$.}
\problemquestion{Find a minimum-weight Steiner tree in $G$.}
\end{problemA}

\begin{problemA}
	\problemtitle{\textsc{Minimum Energy Multicast (MEM)}}
  \problemdefinition{
	 Let $G=(V, E)$ be a graph and let $S\subseteq V$ be a subset of vertices, called terminals. 
	A Steiner tree in $G$ is a tree that spans $S$.		
  }
  \probleminput{A graph $G=(V,E)$, non-negative edge-weights $c$, and $S\subseteq V$.}
\problemquestion{Find a minimum-weight Steiner tree in $G$.}
\end{problemA}


\begin{problemA}
  \problemtitle{\textsc{Minimum Steiner Tree}}
  \problemdefinition{
	 Let $G=(V, E)$ be a graph and let $S\subseteq V$ be a subset of vertices, called terminals. 
	A Steiner tree in $G$ is a tree that spans $S$.		
  }
  \probleminput{A graph $G=(V,E)$, non-negative edge-weights $c$, and $S\subseteq V$.}
\problemquestion{Find a minimum-weight Steiner tree in $G$.}
\end{problemA}


\begin{problemA}
  \problemtitle{\textsc{Range Assignment Problem (RAP)}}
  \problemdefinition{
	 Let $G=(V, E)$ be a graph and let $S\subseteq V$ be a subset of vertices, called terminals. 
	A Steiner tree in $G$ is a tree that spans $S$.		
  }
  \probleminput{A graph $G=(V,E)$, non-negative edge-weights $c$, and $S\subseteq V$.}
\problemquestion{Find a minimum-weight Steiner tree in $G$.}
\end{problemA}

\begin{problemA}
  \problemtitle{\textsc{Minimum Shared Broadcast Tree (SBT)}}
  \problemdefinition{
	 Let $G=(V, E)$ be a graph and let $S\subseteq V$ be a subset of vertices, called terminals. 
	A Steiner tree in $G$ is a tree that spans $S$.		
  }
  \probleminput{A graph $G=(V,E)$, non-negative edge-weights $c$, and $S\subseteq V$.}
\problemquestion{Find a minimum-weight Steiner tree in $G$.}
\end{problemA}

\begin{problemA}
  \problemtitle{\textsc{Minimum Shared Mulicast Tree (SMT)}}
  \problemdefinition{
	 Let $G=(V, E)$ be a graph and let $S\subseteq V$ be a subset of vertices, called terminals. 
	A Steiner tree in $G$ is a tree that spans $S$.		
  }
  \probleminput{A graph $G=(V,E)$, non-negative edge-weights $c$, and $S\subseteq V$.}
\problemquestion{Find a minimum-weight Steiner tree in $G$.}
\end{problemA}

\begin{problemA}
  \problemtitle{\textsc{Symmetric Range Assignment Problem (SRAP)}}
  \problemdefinition{
	 Let $G=(V, E)$ be a graph and let $S\subseteq V$ be a subset of vertices, called terminals. 
	A Steiner tree in $G$ is a tree that spans $S$.		
  }
  \probleminput{A graph $G=(V,E)$, non-negative edge-weights $c$, and $S\subseteq V$.}
\problemquestion{Find a minimum-weight Steiner tree in $G$.}
\end{problemA}

\begin{problemA}
  \problemtitle{\textsc{Weakly Symmetric Range Assignment Problem (WSRAP)}}
  \problemdefinition{
	 Let $G=(V, E)$ be a graph and let $S\subseteq V$ be a subset of vertices, called terminals. 
	A Steiner tree in $G$ is a tree that spans $S$.		
  }
  \probleminput{A graph $G=(V,E)$, non-negative edge-weights $c$, and $S\subseteq V$.}
\problemquestion{Find a minimum-weight Steiner tree in $G$.}
\end{problemA}

             % Appendix A: Something


%--------------------------------------------------------------------------------------------------
% BACK MATTER
%--------------------------------------------------------------------------------------------------
\backmatter

%
% The bibliography
%

%
% Physics-style bibliography
%
%\bibliographystyle{iopart-num}

%
% Modified mathematics bibliography. Uses acm.bst
%
\bibliographystyle{acm}

\bibliography{thesis}

\end{document}

